
\documentclass{article}
\usepackage{colortbl}
\usepackage{makecell}
\usepackage{multirow}
\usepackage{supertabular}

\begin{document}

\newcounter{utterance}

\centering \large Interaction Transcript for game `dond', experiment `semi', episode 1 with nemotron{-}3{-}ultra.
\vspace{24pt}

{ \footnotesize  \setcounter{utterance}{1}
\setlength{\tabcolsep}{0pt}
\begin{supertabular}{c@{$\;$}|p{.15\linewidth}@{}p{.15\linewidth}p{.15\linewidth}p{.15\linewidth}p{.15\linewidth}p{.15\linewidth}}
    \# & $\;$A & \multicolumn{4}{c}{Game Master} & $\;\:$B\\
    \hline

    \theutterance \stepcounter{utterance}  
    & & \multicolumn{4}{p{0.6\linewidth}}{
        \cellcolor[rgb]{0.9,0.9,0.9}{
            \makecell[{{p{\linewidth}}}]{
                \texttt{\tiny{[P1$\langle$GM]}}
                \texttt{Igrate kooperativnu igru pregovaranja u kojoj se trebate dogovoriti s drugim igračem o tome kako podijeliti skup stavki između vas.} \\
\\ 
\texttt{Pravila:} \\
\texttt{(a) Vi i drugi igrač dobivate skup stavki. Svakom od vas također se daje tajna funkcija vrijednosti koja predstavlja koliko cijenite svaku vrstu predmeta.} \\
\texttt{(b) Razmjenjujete poruke s drugim igračem kako biste se dogovorili o tome tko dobiva koje stavke. Svaki može poslati maksimalno 5 poruka ili prekinuti igru ranije dajući tajni prijedlog u bilo kojem trenutku.} \\
\texttt{(c) Od vas se traži da podnese tajni prijedlog koji označava stavke koje želite, formatirane u uglatim zagradama kako slijedi: "{[}Prijedlog: <broj> <naziv stavke>, <broj> <naziv stavke>, <...>{]}"} \\
\texttt{(d) Ako su vaši prijedlozi komplementarni, tj. ima dovoljno stavki za ispunjavanje oba prijedloga, svaki igrač dobiva rezultat na temelju zbroja vrijednosti stavki koje je dobio. U suprotnom, obojica dobivate nula bodova.} \\
\texttt{(e) Vaš cilj je maksimizirati zbroj vašeg rezultata i rezultata drugog igrača.} \\
\\ 
\texttt{Počnimo.} \\
\\ 
\texttt{Skup dostupnih stavki je:} \\
\\ 
\texttt{2 pera, 2 limuni, 2 lonci} \\
\\ 
\texttt{Vaša tajna funkcija vrijednosti je:} \\
\\ 
\texttt{pero: 2, limun: 0, lonac: 3} \\
\\ 
\texttt{VAŽNO: Vaše poruke, osim ako nisu tajni prijedlog, šalju se izravno drugom igraču, pa ne uključujte nikakav odgovor na pravila niti tekst koji najavljuje vašu poruku. Za davanje tajnog prijedloga koristite navedeni format. Ne koristite uglate zagrade pri komuniciranju s drugim igračem ili će se tumačiti kao vaš tajni prijedlog.} \\
            }
        }
    }
    & \\ \\

    \theutterance \stepcounter{utterance}  
    & \multicolumn{4}{p{0.6\linewidth}}{
        \cellcolor[rgb]{0.9,0.9,0.9}{
            \makecell[{{p{\linewidth}}}]{
                \texttt{\tiny{[P1$\rangle$GM]}}
                \texttt{Bok! Vrijedim samo pera (2) i lonce (3), a limuni mi ne vrijede nitišta (0). Ideja mi je da ja uzmem oba pera i oba lanca, a ti bi mogao/la uzeti oba limuna. Kako ti to zvuči?} \\
            }
        }
    }
    & & \\ \\

    \theutterance \stepcounter{utterance}  
    & & & \multicolumn{2}{p{0.3\linewidth}}{
        \cellcolor[rgb]{0.9,0.9,0.9}{
            \makecell[{{p{\linewidth}}}]{
                \texttt{\tiny{[GM$|$GM]}}
                \texttt{continue} \\
            }
        }
    }
    & & \\ \\

    \theutterance \stepcounter{utterance}  
    & & & \multicolumn{2}{p{0.3\linewidth}}{
        \cellcolor[rgb]{0.9,0.9,0.9}{
            \makecell[{{p{\linewidth}}}]{
                \texttt{\tiny{[GM$|$GM]}}
                \texttt{Bok! Vrijedim samo pera (2) i lonce (3), a limuni mi ne vrijede nitišta (0). Ideja mi je da ja uzmem oba pera i oba lanca, a ti bi mogao/la uzeti oba limuna. Kako ti to zvuči?} \\
            }
        }
    }
    & & \\ \\

    \theutterance \stepcounter{utterance}  
    & & \multicolumn{4}{p{0.6\linewidth}}{
        \cellcolor[rgb]{0.9,0.9,0.9}{
            \makecell[{{p{\linewidth}}}]{
                \texttt{\tiny{[GM$\rangle$P2]}}
                \texttt{Igrate kooperativnu igru pregovaranja u kojoj se trebate dogovoriti s drugim igračem o tome kako podijeliti skup stavki između vas.} \\
\\ 
\texttt{Pravila:} \\
\texttt{(a) Vi i drugi igrač dobivate skup stavki. Svakom od vas također se daje tajna funkcija vrijednosti koja predstavlja koliko cijenite svaku vrstu predmeta.} \\
\texttt{(b) Razmjenjujete poruke s drugim igračem kako biste se dogovorili o tome tko dobiva koje stavke. Svaki može poslati maksimalno 5 poruka ili prekinuti igru ranije dajući tajni prijedlog u bilo kojem trenutku.} \\
\texttt{(c) Od vas se traži da podnese tajni prijedlog koji označava stavke koje želite, formatirane u uglatim zagradama kako slijedi: "{[}Prijedlog: <broj> <naziv stavke>, <broj> <naziv stavke>, <...>{]}"} \\
\texttt{(d) Ako su vaši prijedlozi komplementarni, tj. ima dovoljno stavki za ispunjavanje oba prijedloga, svaki igrač dobiva rezultat na temelju zbroja vrijednosti stavki koje je dobio. U suprotnom, obojica dobivate nula bodova.} \\
\texttt{(e) Vaš cilj je maksimizirati zbroj vašeg rezultata i rezultata drugog igrača.} \\
\\ 
\texttt{Počnimo.} \\
\\ 
\texttt{Skup dostupnih stavki je:} \\
\\ 
\texttt{2 pera, 2 limuni, 2 lonci} \\
\\ 
\texttt{Vaša tajna funkcija vrijednosti je:} \\
\\ 
\texttt{pero: 3, limun: 1, lonac: 1} \\
\\ 
\texttt{VAŽNO: Vaše poruke, osim ako nisu tajni prijedlog, šalju se izravno drugom igraču, pa ne uključujte nikakav odgovor na pravila niti tekst koji najavljuje vašu poruku. Za davanje tajnog prijedloga koristite navedeni format. Ne koristite uglate zagrade pri komuniciranju s drugim igračem ili će se tumačiti kao vaš tajni prijedlog.} \\
\\ 
\texttt{Bok! Vrijedim samo pera (2) i lonce (3), a limuni mi ne vrijede nitišta (0). Ideja mi je da ja uzmem oba pera i oba lanca, a ti bi mogao/la uzeti oba limuna. Kako ti to zvuči?} \\
            }
        }
    }
    & \\ \\

    \theutterance \stepcounter{utterance}  
    & & & \multicolumn{4}{p{0.6\linewidth}}{
        \cellcolor[rgb]{0.9,0.9,0.9}{
            \makecell[{{p{\linewidth}}}]{
                \texttt{\tiny{[GM$\langle$P2]}}
                \texttt{Zvuči super! Ja vrijedim pera (3), limune (1) i lonce (1), pa mi taj podijel daje sve što želim (2 limuna = 2 boda) i ostavlja tebi sve što ti najviše vrijedi (2 pera = 6, 2 lonce = 2 → ukupno 8). Za nas двоје to je ukupno 10 bodova, što je maksimum. Slobodno pošalji prijedlog.} \\
            }
        }
    }
     \\ \\

    \theutterance \stepcounter{utterance}  
    & & & \multicolumn{2}{p{0.3\linewidth}}{
        \cellcolor[rgb]{0.9,0.9,0.9}{
            \makecell[{{p{\linewidth}}}]{
                \texttt{\tiny{[GM$|$GM]}}
                \texttt{continue} \\
            }
        }
    }
    & & \\ \\

    \theutterance \stepcounter{utterance}  
    & & & \multicolumn{2}{p{0.3\linewidth}}{
        \cellcolor[rgb]{0.9,0.9,0.9}{
            \makecell[{{p{\linewidth}}}]{
                \texttt{\tiny{[GM$|$GM]}}
                \texttt{Zvuči super! Ja vrijedim pera (3), limune (1) i lonce (1), pa mi taj podijel daje sve što želim (2 limuna = 2 boda) i ostavlja tebi sve što ti najviše vrijedi (2 pera = 6, 2 lonce = 2 → ukupno 8). Za nas двоје to je ukupno 10 bodova, što je maksimum. Slobodno pošalji prijedlog.} \\
            }
        }
    }
    & & \\ \\

    \theutterance \stepcounter{utterance}  
    & & \multicolumn{4}{p{0.6\linewidth}}{
        \cellcolor[rgb]{0.9,0.9,0.9}{
            \makecell[{{p{\linewidth}}}]{
                \texttt{\tiny{[P1$\langle$GM]}}
                \texttt{Zvuči super! Ja vrijedim pera (3), limune (1) i lonce (1), pa mi taj podijel daje sve što želim (2 limuna = 2 boda) i ostavlja tebi sve što ti najviše vrijedi (2 pera = 6, 2 lonce = 2 → ukupno 8). Za nas двоје to je ukupno 10 bodova, što je maksimum. Slobodno pošalji prijedlog.} \\
            }
        }
    }
    & \\ \\

    \theutterance \stepcounter{utterance}  
    & \multicolumn{4}{p{0.6\linewidth}}{
        \cellcolor[rgb]{0.9,0.9,0.9}{
            \makecell[{{p{\linewidth}}}]{
                \texttt{\tiny{[P1$\rangle$GM]}}
                \texttt{{[}Prijedlog: 2 pero, 2 lonac{]}} \\
            }
        }
    }
    & & \\ \\

    \theutterance \stepcounter{utterance}  
    & & & \multicolumn{2}{p{0.3\linewidth}}{
        \cellcolor[rgb]{0.9,0.9,0.9}{
            \makecell[{{p{\linewidth}}}]{
                \texttt{\tiny{[GM$|$GM]}}
                \texttt{proposal} \\
            }
        }
    }
    & & \\ \\

    \theutterance \stepcounter{utterance}  
    & & & \multicolumn{2}{p{0.3\linewidth}}{
        \cellcolor[rgb]{0.9,0.9,0.9}{
            \makecell[{{p{\linewidth}}}]{
                \texttt{\tiny{[GM$|$GM]}}
                [2, 0, 2]
            }
        }
    }
    & & \\ \\

    \theutterance \stepcounter{utterance}  
    & & \multicolumn{4}{p{0.6\linewidth}}{
        \cellcolor[rgb]{0.9,0.9,0.9}{
            \makecell[{{p{\linewidth}}}]{
                \texttt{\tiny{[GM$\rangle$P2]}}
                \texttt{Drugi igrač je poslao tajni prijedlog. Sada je na tebi red da pošalješ vlastiti prijedlog.} \\
            }
        }
    }
    & \\ \\

    \theutterance \stepcounter{utterance}  
    & & & \multicolumn{4}{p{0.6\linewidth}}{
        \cellcolor[rgb]{0.9,0.9,0.9}{
            \makecell[{{p{\linewidth}}}]{
                \texttt{\tiny{[GM$\langle$P2]}}
                \texttt{{[}Prijedlog: 2 pero, 2 lonac{]}} \\
            }
        }
    }
     \\ \\

    \theutterance \stepcounter{utterance}  
    & & & \multicolumn{2}{p{0.3\linewidth}}{
        \cellcolor[rgb]{0.9,0.9,0.9}{
            \makecell[{{p{\linewidth}}}]{
                \texttt{\tiny{[GM$|$GM]}}
                \texttt{proposal} \\
            }
        }
    }
    & & \\ \\

    \theutterance \stepcounter{utterance}  
    & & & \multicolumn{2}{p{0.3\linewidth}}{
        \cellcolor[rgb]{0.9,0.9,0.9}{
            \makecell[{{p{\linewidth}}}]{
                \texttt{\tiny{[GM$|$GM]}}
                [2, 0, 2]
            }
        }
    }
    & & \\ \\

    \theutterance \stepcounter{utterance}  
    & & & \multicolumn{2}{p{0.3\linewidth}}{
        \cellcolor[rgb]{0.9,0.9,0.9}{
            \makecell[{{p{\linewidth}}}]{
                \texttt{\tiny{[GM$|$GM]}}
                [[2, 0, 2], [2, 0, 2]]
            }
        }
    }
    & & \\ \\

\end{supertabular}
}

\end{document}
